% Zone „Das kennst du schon“ – aus bank/tangente-normale-schnittwinkel/zone.jsonl (zusammenbau v0.1)
\einheitenkopf[zone][Kennst du schon]{Das kennst du schon}
\setcounter{aufgabe}{0}

%% TODO Titel: Ich-kann-Satz fehlt in der Bank (Platzhalter: Kettenname) – Nr. 1
%% TODO Anweisung: ein Satz über der ersten Teilaufgabe fehlt in der Bank – Nr. 1
\begin{aufgabe}{Ableitungen bilden}
\begin{teile}
\teil $f(x) = 2x^3 - 5x + 1$ – $f'(x)$? $f'(x) =$ \leerfeld
\teil $f(x) = \frac{1}{3}x^3 - \frac{3}{2}x^2 + 7$ – $f'(x)$? $f'(x) =$ \leerfeld
\end{teile}
\end{aufgabe}

%% TODO Titel: Ich-kann-Satz fehlt in der Bank (Platzhalter: Kettenname) – Nr. 2
%% TODO Anweisung: ein Satz über der ersten Teilaufgabe fehlt in der Bank – Nr. 2
\begin{aufgabe}{Ableitungswert als Tangentensteigung deuten (Grundvorstellung)}
\begin{teile}
\teil Es gilt $f'(2) = 3$. Welche Steigung hat die Tangente an den Graphen von $f$ bei $x = 2$? \leerfeld
\teil Es gilt $f'(-1) = -0{,}4$. Steigt oder fällt die Tangente bei $x = -1$, und welche Steigung hat sie? \leerfeld
\end{teile}
\end{aufgabe}

%% TODO Titel: Ich-kann-Satz fehlt in der Bank (Platzhalter: Kettenname) – Nr. 3
%% TODO Anweisung: ein Satz über der ersten Teilaufgabe fehlt in der Bank – Nr. 3
\begin{aufgabe}{Geradengleichung y = mx + n aus Steigung und Punkt aufstellen, Steigungsdreieck lesen und zeichnen, Nullstelle und y-Achsenabschnitt einer Geraden bestimmen, Punktprobe und Funktionswert}
\begin{teile}
\teil Gerade mit Steigung $2$ durch $P(1 | 5)$ – Gleichung? $y =$ \leerfeld
\teil Gerade mit Steigung $-\frac{1}{2}$ durch $P(4 | 1)$ – Gleichung? $y =$ \leerfeld
\end{teile}
\end{aufgabe}

%% TODO Titel: Ich-kann-Satz fehlt in der Bank (Platzhalter: Kettenname) – Nr. 4
%% TODO Anweisung: ein Satz über der ersten Teilaufgabe fehlt in der Bank – Nr. 4
\begin{aufgabe}{Gleichungen lösen}
\begin{teile}
\teil $3x^2 = 12$ – Lösungen? $x_1 =$ \leerfeld , $x_2 =$ \leerfeld
\end{teile}
\begin{gleichungsraster}[2]
\gl{\text{$x^2 - 2x - 8 = 0$ – Lösungen?}} & \\
\end{gleichungsraster}
\end{aufgabe}

%% TODO Titel: Ich-kann-Satz fehlt in der Bank (Platzhalter: Kettenname) – Nr. 5
%% TODO Anweisung: ein Satz über der ersten Teilaufgabe fehlt in der Bank – Nr. 5
\begin{aufgabe}{Tangens im rechtwinkligen Dreieck und Umkehrfunktion tan⁻¹, Grad-Modus des Rechners}
\begin{teile}
\teil Rechtwinkliges Dreieck: Gegenkathete von $\alpha$ ist $4$ cm, Ankathete $2$ cm – $\mathrm{tan}\alpha$? \leerfeld
\teil $\mathrm{tan}\alpha = 0{,}4$ – $\alpha$ auf eine Stelle? $\alpha \approx$ \leerfeld[°]
\end{teile}
\end{aufgabe}

%% TODO Titel: Ich-kann-Satz fehlt in der Bank (Platzhalter: Kettenname) – Nr. 6
%% TODO Anweisung: ein Satz über der ersten Teilaufgabe fehlt in der Bank – Nr. 6
\begin{aufgabe}{Satz des Pythagoras und Flächenformel des Dreiecks (halbes Produkt der Katheten), Streckenlänge zwischen zwei Punkten}
\begin{teile}
\teil Katheten $3$ cm und $4$ cm – Hypotenuse? \leerfeld[cm]
\teil Länge der Strecke von $A(-1 | 2)$ bis $B(3 | -1)$? \leerfeld
\end{teile}
\end{aufgabe}

%% TODO Titel: Ich-kann-Satz fehlt in der Bank (Platzhalter: Kettenname) – Nr. 7
%% TODO Anweisung: ein Satz über der ersten Teilaufgabe fehlt in der Bank – Nr. 7
\begin{aufgabe}{Ableitungswert als Tangentensteigung deuten (Grundvorstellung) – Fehler finden}
\begin{teile}
\teil Ich finde den Fehler: Für $f(x) = x^2 + 2x$ sagt Mia: „Die Tangente bei $x = 1$ hat die Steigung $f(1) = 3$.“
\end{teile}
\end{aufgabe}

%% TODO Titel: Ich-kann-Satz fehlt in der Bank (Platzhalter: Kettenname) – Nr. 8
%% TODO Anweisung: ein Satz über der ersten Teilaufgabe fehlt in der Bank – Nr. 8
\begin{aufgabe}{Ableitungswert als Tangentensteigung deuten (Grundvorstellung) – selbst rechnen}
\begin{teile}
\teil $f(x) = x^3 - x$ – Steigung der Tangente bei $x = 2$? \leerfeld
\end{teile}
\end{aufgabe}
